\section{Related Work}\label{sec:8}

We situate the main result of SAGA within four groups:
(1) cohomological program analysis; (2) sheaves for architecture and
systems engineering; (3) architecture conformance and formal
connection; (4) mechanized theorem chains and executable measurement.

\subsection{Cohomological program analysis}\label{sec:81}

Young unifies type checking, bug finding, and program equivalence as
\v{C}ech-cohomological analysis of semantic presheaves over observation
sites of Python programs \citep{young2026sheaf}. In its use of local
semantic observation and \v{C}ech $H^1$, and in reporting a Lean
formalization and an executable analyzer, it is the contemporaneous
work closest to SAGA.

The two share the following structure:
\begin{center}
\small
local semantic observations $\to$ site / cover $\to$ presheaf or
coefficient system $\to$ \v{C}ech complex $\to$ $H^1$ obstruction
$\to$ executable analysis
\end{center}

The differences lie in the mathematical center and the level of
abstraction. Young computes cohomology directly on the
program-semantic presheaf of Python. AAT starts from
language-independent Atoms and an equation system, relating facts
spanning different languages, services, and storage representations on
a single architecture object. SAGA constructs semantic-repair
cohomology and equation-generated AAT \v{C}ech cohomology
\emph{separately} and proves a comparison theorem identifying their
$H^1$. The difference is made concrete by the one-cent case: the
amounts in question are all represented as \code{string}, and the
implementation types agree. What SAGA measures is not a difference of
type names but whether locally valid semantic conventions glue into
one global architecture.

\begin{table}[htbp]
\centering
\scriptsize
\begin{tabular}{@{}p{0.20\textwidth}p{0.34\textwidth}p{0.38\textwidth}@{}}
\toprule
Axis & Young 2026 & SAGA \\
\midrule
primary object & Python programs and observation sites &
language-independent Atom families and architecture objects \\
coefficient & semantic presheaf, $\Ftwo$ realization & the
$\QE=\ObsE/\IOb$ generated by the equation system \\
central theorem & program-analysis claims via \v{C}ech cohomology &
the comparison $H^1_{\mathrm{sem}} \cong \check H^1(\cU,\QE)$ \\
repair reading & rank = independent fixes & residual classes,
boundary, repair gluing \\
empirical unit & 375 program-analysis benchmarks & a real microservice
architecture with before/after-repair comparison \\
domain-specific construction & direct computation on the
program-semantic presheaf & two independent presentations (semantic
repair / equation quotient), the construction of $\chi/\Phi/\kappa$,
exactness conditions relativized to architecture data \\
discharged obligation & connecting cohomology computations to
program-analysis claims & derivation of soundness, checking of the
$\ker$/$\operatorname{im}$ conditions, proof of the semantic
correspondence of residual classes \\
\bottomrule
\end{tabular}
\caption{Comparison of Young 2026 and SAGA.}
\label{tab:young}
\end{table}

What SAGA claims as the novelty of the comparison is the last two rows
of Table~\ref{tab:young} --- the domain-specific construction and the proof
obligations discharged there --- not the general mechanism by which a
coefficient isomorphism induces a cochain isomorphism (the attribution
of what is proper to SAGA is in \S\ref{sec:51}).

Young's Lean formalization (1{,}259 lines), the 375 benchmarks, and
the evaluation numbers are cited as results reported by that paper.
Since no public link to the code artifact of that paper's analyzer
could be confirmed, source-level comparison is outside the scope of
this paper.

In the lineage of global-section obstructions, we place
sheaf-theoretic contextuality \citep{abramsky2011sheaf} and its
non-vanishing obstructions in \v{C}ech cohomology
\citep{abramsky2012cohomology} as foundational references. As prior
examples of computable obstructions, we cite the application of
\v{C}ech cohomology to CSP and structure isomorphism
\citep{oconghaile2022cohomology}, and the task sheaf for distributed
task solvability \citep{felber2025sheaf}.

\subsection{Sheaves for architecture and systems engineering}\label{sec:82}

Gibson gives a Lean-verified sheaf model for multi-view consistency in
model-based systems engineering, characterizing global designs by
compatible local designs on pairwise interfaces
\citep{gibson2026sheaves}. SAGA develops the complementary
obstruction-theoretic direction: constructing the $H^1$ class that
remains when local data do not glue, and proving the comparison
between its semantic realization and its equation-generated
realization.

As the historical starting point we place the sheaf semantics of
objects and interaction \citep{goguen1992sheaf}. As foundational
references for cellular sheaves and finite computation we cite
\citet{curry2014sheaves}, \citet{robinson2017sheaves}, and
\citet{hansen2019spectral}.

\subsection{Architecture conformance and formal connection}\label{sec:83}

The architectural-mismatch analysis of \citet{garlan1995architectural}
examined the mismatch of assumptions embedded in components,
connectors, and the construction process. SAGA moves this kind of
mismatch into the mathematics of local equations and global gluing,
treating the residue as an $H^1$ class. Reflexion models \citep{murphy1995reflexion}, which
compute convergence / divergence between a high-level model and a
source model, are the representative starting point of conformance.
Wright \citep{allen1997formal} formalized the compatibility of
connector protocols in CSP. What SAGA treats is not protocol
conformance but the cohomology class of whether multiple local
semantic repairs glue across a cover. Compared with the survey of the
expressive power of ADLs \citep{medvidovic2000classification} and the
systematic surveys of architecture erosion
\citep{desilva2012controlling, li2022understanding}, SAGA offers a
global obstruction class for selected local data and a repair
comparison.

\subsection{Mechanized theorem chains and executable measurement}\label{sec:84}

Lean~4 \citep{demoura2021lean} and Mathlib \citep{mathlib2020lean} are
cited as the provenance of the formalization environment. The
formalization contribution of SAGA is not the use of Lean itself, but
the construction, as a single machine-checked theorem chain reaching
the conclusion bundle of the central theorem, of the semantic repair
presentation, the cover-relative \v{C}ech complex, the coefficient
isomorphism and cochain comparison, true sheaf descent, and the finite
witnesses of the zero and nonzero cases (Section~\ref{sec:6}). The
comparison with the formalizations of Young and Gibson is made in
terms of the scope of the theorem chain and the connection to
executable measurement.

\subsection{Synthesis}\label{sec:85}

Prior work has established sheaves as a language for global
consistency, cohomology as computable obstruction, and formal
architecture models as a basis for conformance analysis. To these
lineages SAGA adds the construction of an equation-generated AAT
\v{C}ech complex for software architecture and the proof that its
first cohomology agrees with an independently defined semantic-repair
obstruction. Lean verifies the comparison, and ArchSig evaluates its
finite architectural instances.

\textbf{TODO:} at submission time, re-verify the latest versions and
publication status of the 2026 references.
